\section{Related Work}
\label{sec:rw}

\noindent\textbf{Datasets.} Digital entertainment such as movies and TV dramas are attracting much research attention from the multimodal research field, as they offer a valuable snapshot of situations and scenarios in the real-world. 
LSMDC~16~\citep{lsmdc2015cvpr,lsmdc2017ijcv}, a combination of sub-datasets in several related domains, is the pioneer in this direction. 
Its success sparked the creation of a number of follow-ups: Pororo~\cite{pororo2017arxiv} which contains animation internal dialogues and external human summaries, YMS~\citep{yms2018cvpr} which comprises human-narrated movie plot summaries, and CMD~\citep{cmd2020accv} which pairs key clips with descriptions. Later resources including MovieNet~\citep{movienet2020eccv}, MAD~\citep{mad2022cvpr}, and SYMON~\citep{symon2022arxiv} not only expanded the scale of data and added more detailed descriptions (see statistics in Tab.~\ref{tab:compare}), but also introduced more diverse tasks, e.g., genre prediction, segment ordering, and multimodal retrieval. 

However, as discussed in~\cref{sec:introduction}, these corpora are limited in three major aspects (i.e., segmentation according to shots/scenes rather than a slot, lack of empathy and context in the narratives, and absence of long-form relationships due to short duration), constraining the scope and utility. The proposal of Movie101~\citep{yue2023movie101} gives a feasible solution, but its data cannot be obtained directly, and it lacks detailed plot descriptions. The \TVD dataset, in contrast, is constructed to mitigate these issues.

\noindent\textbf{Models.} Although recent years have witnessed many studies that conducted experiments using the above cinema/TV multimodal datasets~\citep{lsmdc2017ijcv,movienet2020eccv,symon2022arxiv,movie22022arxiv,yu2021transitional}, to the best of our knowledge, the model weights of \textit{neither} the methods proposed in these works \textit{nor} the baselines introduced in the original corpus papers are publicly available. Moreover, the architectures of these algorithms vary as they focus on different problems. As a result, it is difficult for follow-up works in this research vein (including ours) to recycle the models built in earlier studies, leading to extra training costs, unnecessary carbon footprint, as well as barriers to comparing different algorithms in a fair way.

To address this issue, for the first time, we align the level of machine learning tasks in the cinema/TV domain with the stage of human intelligence defined by Cognitive Scientists. Aiming to provide a strong and reusable baseline, we proposed a framework that can solve different tasks with a unified network design. We promise to release our model weights, so that follow-up works no longer need to develop their approaches from scratch.

\noindent\textbf{Bullet Screen Comments.} Bullet Screen Comments (BSCs) refer to the fashionable type of comments that fly across the screen like bullets when the video is playing. On the video timeline, a BSC appears at the point where it is initially posted by the user and hovers for a few seconds.
Compared with traditional static comments, BSCs are more likely to be strongly associated with specific dramatic pieces.
They make video viewers feel like having ``real-time'' discussions with others. Users can also make active interactions by upvoting BSCs they favour. 
% See~\cref{sec:app_bsc} for a real-world screenshot.
See the appendix for a real-world screenshot.

A group of studies has validated the usefulness of BSCs in helping multimodal methods to model video semantics~\citep{bs32018movie,bs2www2020,bs12022aaai,yang2019time,wang2020videoic}. However, as far as we are aware, none of these works have the copyright of the BSCs, and there is no \textit{direct} download URLs have been provided. Besides, they have never investigated the role of BSCs in processing dramatic plots, either. In \TVD, we included more than 26 million BSC sentences, making the first public BSC dataset. Additionally, it is also the largest text corpus in the cinema/TV domain to date, whose scale is sufficient to bootstrap the training of large language models. Moreover, for the first time, we empirically justified that BSCs can effectively enhance systems' ability to modelling complex plots. 